\chapter{Range-Extended Network-Aware Model Partitioning Framework} \label{chap:ch4}

\section{Approach} \label{sec:se41}

After analyzing the state of the art on model partitioning and \ac{ML} 
orchestration in edge computing, we found a lack of adaptation of models
with cyclic topology (e.g., \acp{RNN}, MobileNets, etc.) to the edge
environment. Additionally, current frameworks tend to make use of less
constrained devices and use more powerful and even less constrained
devices in the later stages of the inference tasks, which may not be the
most optimal approach in a network where the devices lay closer together.

To address these issues, we propose a novel model partitioning framework that
exploits the cyclic topology of certain \ac{ML} models to use devices as
loop processors, reusing them in different stages of the inference task.
The presented approach will mostly make use of highly constrained devices with a
lower usage of less constrained devices. Paired with including devices with more 
computing power, this approach allows us to better distribute the computational 
load across the available devices, taking advantage of their proximity in the 
network.

\section{Methods} \label{sec:se42}

For the testing and validation of the proposed framework, we will use a set of
15 \ac{MCU}-based devices with available memory on the order of kilobytes and one
\ac{MPU}-based device with available memory on the order of megabytes. 

The devices will be connected via Wi-Fi and deployed in an 
indoor environment. The network topology represents a typical smart building 
scenario, with the network representing a fully-connected graph of 16 nodes.

The metrics for evaluation will include inference latency, energy consumption, and
memory usage per device. We will then compare the performance of our proposed
framework against the state of the art model partitioning frameworks that we
analysed in Chapter \ref{chap:ch3}.

\section{Expected Results} \label{sec:se43}

We expect to achieve a successful deployment of partitioned models with cyclic 
components, extending the range of \ac{DNN} models that can be deployed in edge
environments. 

Additionally, we anticipate improved performance in terms of
latency and energy efficiency compared to existing model partitioning frameworks,
especially in scenarios with a high density of constrained devices. This is 
partially due  to improved communication methods between devices and better 
utilization of edge devices among the far-edge devices in the network.

\section{Project plan} \label{sec:se44}

The project is planned to be executed over a period of 5 months, resulting in 
18 available weeks, divided into the following tasks:

\begin{description}[style=unboxed, leftmargin=1cm, font=\bfseries]

    \item[Testbed Setup (Weeks 1-8):]
    in this task, we will set up the hardware and software environment for testing,
    including the configuration of the devices and network. Furthermore, we will
    implement the necessary tools for monitoring and evaluating performance metrics.
    
    \item[Deployment of the Model Partition Segments (Weeks 4-7):]
    this task involves the deployment of the partitioned model segments onto the
    selected devices, ensuring correct functionality and efficient communication
    between them.
    
    \item[Partitioning Subprocedure Implementation (Weeks 7-12):]
    responsible for implementing the core partitioning algorithm that will divide
    the \ac{DNN} models into lightweight segments suitable for the constrained 
    devices in the network. Ideally, this subprocedure will take into account the 
    cyclic topology of certain models to optimize the partitioning process.
    
    \item[Testing and Evaluation (Weeks 12-14):]
    during these weeks, we will conduct extensive testing of the deployed framework,
    measuring the defined performance metrics and comparing them against baseline
    results from existing frameworks.

    \item[Dissertation Elaboration (Weeks 15-18):]
    after completing the implementation work and testing, we will focus on writing
    the dissertation, documenting the methodology, results, and conclusions of the
    project.

    
\end{description}

The first three tasks represent the milestone corresponding to the MVP of the
project, expected to be completed by the end of week 12. The testing and evaluation
task will provide the necessary data to analyze the performance of the proposed
framework, while the dissertation elaboration task will ensure that the findings
are well-documented and presented. 

\section{Conclusions} \label{sec:se45}

This chapter proposes a novel framework for closely integrating constrained devices
to varying degrees, enabling the deployment of \ac{DNN} models with cyclic 
topology in edge environments. This marks a significant advancement in the
field of edge computing and \ac{IoT}, opening new possibilities for applications
that require real-time processing and low latency.
